\documentclass[11pt, a4paper, oneside]{article}
\usepackage[utf8]{inputenc}
\usepackage[margin=2cm]{geometry} 
\usepackage[usenames,dvipsnames]{color} 
\usepackage{fullpage}
\usepackage{hyperref} 
\usepackage{uri}
\hypersetup{colorlinks, citecolor=Violet, linkcolor=Red, urlcolor=Blue}
\urlstyle{same}
%\renewcommand{\baselinestretch}{1.2}
\usepackage{graphicx} 
\usepackage{amssymb} 
\usepackage{libertine}
%\usepackage{showframe}
\DeclareGraphicsRule{.tif}{png}{.png}{`convert #1 `dirname
  #1`/`basename #1 .tif`.png}
%\usepackage{sectsty}
%\allsectionsfont{\sffamily\mdseries\upshape} 
\usepackage{tabularx}

\makeatletter
\g@addto@macro{\UrlBreaks}{\UrlOrds}
\makeatother

% sans serifs for whole document.
% \renewcommand\familydefault{\sfdefault}

\begin{document}
\thispagestyle{empty}

\noindent
\begin{tabularx}{\textwidth}{X  r}

\huge\textbf{Dr Charles Patrick Martin} & \\ 
\vspace{0.2cm}&\\
School of Computing & P: +61 2 6125 3139 \\  
Hanna Neumann Building 145, Science Road & M: +61 419 176 858\\ 
The Australian National University & charles.martin@anu.edu.au\\ 
Canberra ACT 2601 & \url{https://charlesmartin.au}\\
\end{tabularx}

\noindent
\begin{tabularx}{\textwidth}{l  p{8cm}}
   \raisebox{-\totalheight}{\includegraphics[width=0.4\textwidth]{charlesmartin-synth-profile}} &

  \vspace{0.01cm}
  \large \textbf{Charles Martin} is a computer scientist specialising in music
  technology, musical AI, and human computer interaction. He is also a
  performer, creator, and teacher.

\end{tabularx}

\subsection*{Academic Positions}
\label{researchexperience}

\begin{itemize}
	\item \textbf{2022 ---} Senior Lecturer, School of Computing, The Australian National University
	\item \textbf{2019 --- 2021} Lecturer, School of Computing, The Australian National University
\item \textbf{2016 --- 2019} Postdoctoral Fellow, University of Oslo,
  Department of Informatics, Robotics and Intelligent Systems Group
  and RITMO Centre for Interdisciplinary Studies in Rhythm, Time and
  Motion
% \item \textbf{2013 --- 2016} PhD Student, The Australian National
%   University, Research School of Computer Science. 
\end{itemize}

\subsection*{Qualifications}
\begin{itemize}
	\item Senior Fellow of the Higher Education Academy (SFHEA)
	\item Senior Member of the Association for Computing Machinery (ACM)
	\item IEEE Senior Member
\end{itemize}

\subsection*{Education}
\label{education}

\begin{itemize}

\item Doctor of Philosophy (Computer Science).\\
  \textbf{The Australian National University} 2016 Canberra, Australia.
\item Master of Fine Arts (Musical Performance).\\
  \textbf{Lule{\aa} University of Technology} 2012 Pite{\aa}, Sweden.
\item Master of Music (Percussion and Computer Music).\\
  \textbf{The Australian National University} 2009.
\item Bachelor of Science with First Class Honours (Mathematics).\\
  \textbf{The Australian National University} 2008.
\end{itemize}

\subsection*{Honorary and Visiting Positions}

\begin{itemize}
\item \textbf{2021 ---} Visiting Scientist, CSIRO
\item \textbf{2021 --- 2022} Associate Professor (F{\o}rsteamanuensis), University of Oslo,
  Department of Informatics
\end{itemize}

\subsection*{Career Breaks}
\begin{itemize}
	\item \textbf{2021 Jan--Jul} Parental Leave
	\item \textbf{2023 Jul--Jan} Parental Leave
\end{itemize}

\subsection*{Service and Leadership}
\label{serviceleadership}

\begin{itemize}

\item General Chair of International Conference on New Interfaces for Musical Expression (NIME) 2025 (hosted at ANU)

\item Assistant Associate Director (Education), ANU School of Computing (2025 ---)

\item Deputy Sub-Dean (Education), ANU College of Engineering, Computing and Cybernetics (2024)

\item Committee Member on the College of Systems and Society Delegated Ethics	Review Committee (CSS DERC) 2024---2026
  
\item General Chair of OzCHI 2022 conference (hosted at ANU)

\item Co-Chair of ``Generative AI and HCI Workshop" at CHI 2022---2025

\item President (2019 --- 2021), Secretary (2021 ---) of the Australasian
	Computer Music Association, Australia's primary academic community for
	computer and electronic music.

\item Lead editor for \emph{Chroma: Journal of the Australasian Computer Music
	Association}. \url{https://journal.computermusic.org.au} (2021---2023)

\item Early Career Academic Fellow of University House (2020)

\item Paper and Artist Talk Chair of the Australasian Computer Music Conference (online/ANU)
	in 2020.

\item NIME Board Member (2024--), Web Officer (2019 --- 2024)
  
\item NIME workshop chair in 2022, demo chair in 2016.

\item Organised NIME workshops: ``Building NIMEs with Embedded AI'' (2024),
	``Critical Perspectives on AI/ML in Musical Interfaces'' (2020, 2021)
	and ``Making Predictive NIMEs with Neural Networks'' (2019).

\item AC (associate chair) for ACM DIS 2022, ACM MM 2024, ACM CHI 2026. Special reviewer recognition for CHI 2018 and DIS 2022.

\item reviewer for ERC Consolidator grant (2025)

\end{itemize}


\subsection*{Research Mentorship}
\label{research-mentorship}

\subsubsection*{PhD Supervision (current)}

\begin{itemize}

\item Xinlei Niu: funded by ANU Scholarship, 2022-- (submitted, awaiting examination)

\item Minsik Choi: ANU Scholarship, 2025--

\item Sandy Ma: funded by CSIRO Next Generation Graduate Program, 2025--

\item Albert Ngabo-Niyonsenga: funded by Australian Government Research Training Program (RTP), 2025--

\end{itemize}

\subsubsection*{PhD Supervision (completed)}

\begin{itemize}

\item Yichen Wang: "Beyond Mapping: Exploring Spatial Musical Interface Design with Augmented Reality", ANU (accepted, awaiting conferral) (2021--2026)

\item Dr Benedikte Wallace: "", University of Oslo (2018--2023)

\item Dr Tønnes Nygaard: "", University of Oslo, co-supervision with Jim Tørresen, Kyrre Glette (2016--2020)

\end{itemize}

\subsubsection*{Master/honours Supervision}

\begin{itemize}

\item 14+ Master students supervised (including MSc by research, Master by coursework).

\item 16+ Honours (undergraduate) students supervised.

\end{itemize}

\subsubsection*{PhD Examinations}

2 in progress. Also 1 MA by Research examination.

\subsection*{Publications and Research Outputs}
\label{publications}

Google Scholar page: \url{https://scholar.google.com/citations?user=mTlH4G8AAAAJ}

\subsubsection*{Conference Papers, Journal Articles, and Book Chapters}

\begin{enumerate}

\item SteerMusic: Enhanced Musical Consistency for Zero-shot Text-Guided and Personalized Music Editing (X Niu, KW Cheuk, J Zhang, N Murata, C-H Lai, M Mancusi, W Choi, G Fabbro, W-H Liao, CP Martin, Y Mitsufuji). \emph{AAAI Conference on Artificial Intelligence}, 2026. \arxiv{2504.10826}

\item Beyond Video-to-SFX: Video to Audio Synthesis with Environmentally Aware Speech (X Niu, J Ma, D Harper-Harris, X Zhang, CP Martin, J Zhang). \emph{ICASSP}, 2026. \arxiv{2509.15492}

\item Exploring the Impact of Spatial Awareness on Large-Scale AR DMIs (Q Zhou, CP Martin). \emph{International Conference on New Interfaces for Musical Expression}, 2025.

\item Luna: An AR Musical Instrument on the Meta Quest 2 (S Dietz, CP Martin). \emph{International Conference on New Interfaces for Musical Expression}, 2025.

\item Touching Wires: Tactility and a Quilted Musical Interface for Human-AI Musical Co-Creation (S Ma, CP Martin). \emph{International Conference on New Interfaces for Musical Expression}, 2025.

\item Seeing the Sound: Supporting Musical Collaboration with Augmented Reality (Y Wang, M Xi, M Adcock, CP Martin). \emph{C\&C '25: Proceedings of Creativity and Cognition}, 2025. \doi{10.1145/3698061.3726905}

\item AI See, You See: Human-AI Musical Collaboration in Augmented Reality (Y Wang, CP Martin). \emph{CHI EA '25: Proceedings of the Extended Abstracts of the CHI Conference on Human Factors in Computing Systems}, 2025. \doi{10.1145/3706599.3720052}

\item A Musical Tension and Release Framework for Sonic Interaction Designers (M Choi, J Andres, A Hunter, CP Martin). In \emph{CHI EA '25:	Proceedings of the Extended Abstracts of the CHI Conference on Human Factors in Computing Systems}, 2025. \doi{10.1145/3706599.3720089}

\item Scenario-Based Design to Envision How GenAI Can Support Sound Design Practices (M Choi, J Andres, A Hunter, CP Martin). \emph{Proceedings of OzCHI}, 2024.

\item Tonal Cognition in Sonification: Exploring the Needs of Practitioners in Sonic Interaction Design (M Choi, J Andres, CP Martin) In \emph{Proceedings of the 19th International Audio Mostly Conference: Explorations in Sonic Cultures}, 2024. \doi{10.1145/3678299.3678321}

\item HybridVC: Efficient Voice Style Conversion with Text and Audio Prompts (X Niu, J Zhang, CP Martin) In \emph{Proceedings of Interspeech 2024}, 2024. \arxiv{2404.15637}

\item Latent Optimal Paths by Gumbel Propagation for Variational Bayesian Dynamic Programming (X Niu, C Walder, J Zhang, CP Martin) In \emph{Proceedings of the 41st International Conference on Machine Learning}, 2024. \arxiv{2306.02568}

\item SoundLoCD: An Efficient Conditional Discrete Contrastive Latent Diffusion Model for Text-to-Sound Generation (X Niu, J Zhang, C Walder, CP Martin) In \emph{IEEE International Conference on Acoustics, Speech and Signal Processing (ICASSP)}, 2024. \doi{10.1109/ICASSP48485.2024.10446349}

\item Breaking from Realism: Exploring the Potential of Glitch in AI-Generated Dance (B Wallace, K Nymoen, J Torresen, CPMartin) \emph{Digital Creativity}. 2024. \doi{10.1080/14626268.2024.2327006}

\item Composing Interface Connections for a Networked Touchscreen Ensemble (CP Martin, A Hunter, B Schuetze, Y Wang). In \emph{Proceedings of the International Workshop on the Internet of Sounds}, 2023. \doi{10.1109/IEEECONF59510.2023.10335226}

\item Mobility, Space and Sound Activate Expressive Musical
Experience in Augmented Reality (Y Wang, M Xi, M Adcock, CP Martin). In \emph{Proceedings of the
International Conference on New Interfaces for Musical
Expression}, 2023. \url{https://www.nime.org/proceedings/2023/nime2023_17.pdf}

\item Embodying an Interactive AI for Dance Through Movement Ideation 
(B Wallace, C Hilton, K Nymoen, J Torresen, CP Martin, R Fiebrink). 
In \emph{Creativity and Cognition}, 2023. \doi{10.1145/3591196.3593336}

\item GenAICHI 2023: Generative AI and HCI at CHI 2023 (M Muller, LB Chilton, A
Kantosalo, QV Liao, ML Maher, CP Martin, and G Walsh). In \emph{CHI Conference on Human
Factors in Computing Systems Extended Abstracts (CHI EA '23)}, 2023. \doi{10.1145/3544549.3573794}

\item Cubing Sound: Designing a NIME for Head-mounted Augmented Reality (Y Wang, C Martin). In  \emph{Proceedings of the
    International Conference on New Interfaces for Musical
    Expression}, 2022. \doi{10.21428/92fbeb44.b540aa59}
    
\item Spatial-Temporal-Class Attention Network for Acoustic Scene Classification 
(X Niu, CP Martin). In \emph{IEEE International Conference on Multimedia and Expo}, 
2022. \doi{10.1109/ICME52920.2022.9859735}

\item GenAICHI: Generative AI and HCI (M Muller, LB Chilton, A
		Kantosalo, CP Martin, and G Walsh). In \emph{CHI Conference on Human
		Factors in Computing Systems Extended Abstracts (CHI EA '22)}, 2022.
		\doi{10.1145/3491101.3503719}
		
\item Augmenting Sculpture with Immersive Sonification (Y Wang, H Gardner, C Martin, M Adcock). In \emph{IEEE Conference on Virtual Reality and 3D User Interfaces Abstracts and Workshops (VRW)}, 2022. \doi{10.1109/VRW55335.2022.00164}

\item Designing Sound Synthesis Interfaces for Head-mounted Augmented Reality (Y Wang, C Martin). In \emph{IEEE Conference on Virtual Reality and 3D User Interfaces Abstracts and Workshops (VRW)}, 2022. \doi{10.1109/VRW55335.2022.00078}

\item Performing with a Generative Electronic Music Controller (CP Martin). In \emph{Joint Proceedings of the ACM IUI Workshops}, 2022. \url{https://ceur-ws.org/Vol-3124/paper10.pdf}

\item Real-world Embodied AI Through a Morphologically Adaptive
  Quadruped Robot (TF Nygaard, CP Martin, J Torresen, K
  Glette, D Howard) \emph{Nature Machine Intelligence}. 2021.
  \doi{10.1038/s42256-021-00320-3}

\item Environmental Adaptation of Robot Morphology and Control Through
  Real-world Evolution (TF Nygaard, CP Martin, D Howard, J
  Torresen, K Glette) \emph{Evolutionary Computation}.
  2021. \arxiv{2003.13254} \doi{10.1162/evco_a_00291}


\item Learning Embodied Sound-Motion Mappings: Evaluating AI-Generated Dance
	Improvisation (B Wallace, CP Martin, J Torresen, and K Nymoen).
	\emph{Creativity and Cognition}, 2021. \doi{10.1145/3450741.3465245}

\item Exploring the Effect of Sampling Strategy on Movement Generation
  with Generative Neural Networks (B Wallace, CP Martin, J
  Torresen, and K Nymoen). \emph{Proceedings of EvoMusArt}, 2021.
  \doi{10.1007/978-3-030-72914-1_23}

\item An Application for Evolutionary Music Composition using
  Autoencoders (RN McArthur and CP Martin). \emph{Proceedings of
    EvoMusArt}, 2021. \doi{10.1007/978-3-030-72914-1_29}


\item Data Driven Analysis of Tiny Touchscreen Performance with
  MicroJam (CP Martin and J Torresen) \emph{Computer Music
    Journal}. 2020. \doi{10.1162/COMJ_a_00536}

\item Understanding Musical Predictions with an Embodied Interface for
  Musical Machine Learning (CP Martin, K Glette, TF Nygaard, and J
  Torresen) \emph{Frontiers in Artificial Intelligence}. 2020.
  \doi{10.3389/frai.2020.00006}

\item Towards Movement Generation with Audio Features (B Wallace, CP Martin, J Torresen, and K Nymoen). \emph{Proceedings of the
    11th International Conference on Computational Creativity}, 2020.
  \arxiv{2011.13453}

\item Comparing three data representations for music with a
  sequence-to-sequence model (S Li, CP Martin). \emph{Australasian
    Joint Conference on Artificial Intelligence}, 2020.
  \doi{10.1007/978-3-030-64984-5_2}

\item A laptop ensemble performance system using recurrent neural
  networks (R Proctor and CP Martin). \emph{Proceedings of the
    International Conference on New Interfaces for Musical
Expression}, 2020. \doi{10.5281/zenodo.4813481}

\item Sonic sculpture: Activating engagement with head-mounted
  augmented reality (CP Martin, Z Liu, Y Wang, W He, and H
  Gardner). \emph{Proceedings of the International Conference on New
  Interfaces for Musical Expression}. 2020. \doi{10.5281/zenodo.4813445}

\item Towards playing in the `air': Modeling motion-sound energy
  relationships in electric guitar performance using deep neural
  networks (C Erdem, Q Lan, J Fuhrer, CP Martin, J Torresen,
  and AR Jensenius). \emph{Proceedings of the 17th Sound and Music
    Computing Conference}. 2020. \nbn{URN:NBN:no-82501}
  
\item An Interactive Musical Prediction System with Mixture Density
  Recurrent Neural Networks (CP Martin, J Torresen).
  \emph{Proceedings of the International Conference on New Interfaces
  for Musical Expression}. 2019. \doi{10.5281/zenodo.3672952} 

\item Generating Convincing Harmony Parts with Bidirectional Long
  Short-Term Memory Networks (A Faitas, S Baumann, T R Næss, J
  Torresen, CP Martin). \emph{Proceedings of the International
  Conference on New Interfaces for Musical Expression}. 2019. \doi{10.5281/zenodo.3672980}
 
\item A Physical Intelligent Instrument using Recurrent Neural
  Networks (TR Næss and CP Martin). \emph{Proceedings of the
    International Conference on New Interfaces for Musical
Expression}. 2019. \doi{10.5281/zenodo.3672874}

\item Free-improvised rehearsal-as-research for musical HCI (CP
  Martin, H Gardner). In T Mudd, S Holland, K Wilkie, A
  McPherson, and M Wanderley (Eds.), \emph{New Directions in Music
    and Human-Computer Interaction}, Springer. 2019.
  \doi{10.1007/978-3-319-92069-6_17}
  
\item Self-Modifying Morphology Experiments with DyRET: Dynamic Robot
  for Embodied Testing (TF Nygaard, CP Martin, J
  Tørresen, K Glette). \emph{Proc. of the IEEE Int. Conf. on Robotics
  \& Automation (ICRA)}. 2019. \doi{10.1109/ICRA.2019.8793663} 

\item Comparing Models for Harmony Prediction in an Interactive Audio
  Looper (B Wallace and CP Martin). \emph{International Conference on
    Computational Intelligence in Music, Sound, Art and Design}. 2019.
  \doi{10.1007/978-3-030-16667-0_12}

\item Identifying Reusable Early-Life Options (A Weber, CP Martin,
  J Torresen, and BC da Silva).  \emph{International
    Conference on Development, Learning and Epigenetic Robotics
    (ICDL-EpiRob)}. 2019. \doi{10.1109/DEVLRN.2019.8850725}

\item Evolving Robots on Easy Mode: Towards a Variable Complexity
  Controller for Quadrupeds (TF Nygaard, CP Martin, J Torresen, K
  Glette). \emph{International Conference on the Applications of
    Evolutionary Computation}. 2019.
  \doi{10.1007/978-3-030-16692-2_41}

\item Experiences from real-world evolution with dyret: Dynamic robot
  for embodied testing (TF Nygaard, J Nordmoen, KO Ellefsen, CP Martin, J Tørresen, and K Glette). \emph{Nordic Artificial
    Intelligence Research and Development}. 2019.
  \doi{10.1007/978-3-030-35664-4_6}

\item Real-World Evolution Adapts Robot Morphology and Control to
  Hardware Limitations (TF Nygaard, CP Martin, E Samuelsen, J
  Torresen, and K Glette). \emph{Proceedings of the Genetic and
    Evolutionary Computation Conference}. 2018.
  \doi{10.1145/3205455.3205567} 
  
\item Composing an Ensemble Standstill Work for Myo and Bela (CP
  Martin, AR Jensenius, and J Torresen) \emph{Proceedings of the
    International Conference on New Interfaces for Musical
    Expression}. 2018. \doi{10.5281/zenodo.1302543}

\item Bela-Based Augmented Acoustic Guitars for Inverse Sonic
  Microinteraction (VEG Sanchez, A Zelechowska,
  CP Martin, KAV Bjerkestrand, V Johnson, and AR
  Jensenius).  \emph{Proceedings of the
    International Conference on New Interfaces for Musical
    Expression}.  2018. \doi{10.5281/zenodo.1302599}

\item Exploring Mechanically Self-Reconfiguring Robots for Autonomous
  Design (TF Nygaard, CP Martin, J Torresen, K
  Glette). In \emph{ICRA18 Workshop on Autonomous Robot Design}. 2018.
  \arxiv{1805.02965}

\item RoboJam: A Musical Mixture Density Network for Collaborative
  Touchscreen Interaction. (CP Martin and J Torresen).
  \emph{Evolutionary and Biologically Inspired Music, Sound, Art and
    Design: Proceedings of EvoMUSArt '18}. 2018.
  \doi{10.1007/978-3-319-77583-8_11} \arxiv{1711.10746}


\item Percussionist-Centred Design for Touchscreen Digital Musical
  Instruments. (CP Martin). \emph{Contemporary Music Review}. 2017.
  \doi{10.1080/07494467.2017.1370794}

\item Deep Models for Ensemble Touch-Screen Improvisation.
  (CP Martin, KO Ellefsen, and J Torresen). In
  \emph{Proceedings of Audio Mostly}. 2017.
  \doi{10.1145/3123514.3123556}

\item Exploring social mobile music with tiny touch-screen
    performances. (CP Martin and J Torresen).  In
  \emph{Proceedings of the 14th Sound and Music Computing Conference
    (SMC)}. 2017. \hdl{10852/57825}

\item Microjam: An app for sharing tiny touch-screen
    performances. (CP Martin, J Torresen). In \emph{Proceedings
    of the International Conference on New Interfaces for Musical
    Expression}. 2017. \hdl{10852/56065}

\item Can Machine Learning Apply to Musical Ensembles? (C Martin and
  H Gardner) In \emph{Proceedings of the CHI Human-Centered Machine
    Learning Workshop}. 2016. \doi{10.5281/zenodo.56379}

\item Free-Improvised Rehearsal-as-Research for Musical HCI. In
  \emph{Proceedings of the CHI Musical HCI Workshop}. 2016.
  \doi{10.5281/zenodo.56378}

\item Intelligent agents and networked buttons improve
    free-improvised ensemble music-making on touch-screens (C
  Martin, H Gardner, B Swift, \& M Martin). In \emph{Proceedings of
    the SIGCHI Conference on Human Factors in Computing System (CHI)}.
    2016. \doi{10.1145/2858036.2858269}

\item A percussion-focussed approach to preserving touch-screen
    improvisation (C Martin, H Gardner). In D England, N
  Bryan-Kinns, \& T Schiphorst (Eds.), \emph{Curating the Digital: Spaces
  for Art and Interaction}, Springer Series on Cultural Computing.
  Switzerland, Springer International Publishing. 2016.
  \doi{10.1007/978-3-319-28722-5_5}

\item Music of 18 performances: Evaluating apps and agents with
    free improvisation (C. Martin, H. Gardner, B. Swift, \& M.
  Martin). In \emph{Proceedings of the Australasian Computer Music
    Conference (ACMC)}. 2015. \hdl{1885/95205}

\item Tracking Ensemble Performance on Touch-Screens with Gesture
  Classification and Transition Matrices (C. Martin, H. Gardner, \& B.
  Swift). In \emph{Proceedings of the International Conference on New
    Interfaces for Musical Expression (NIME)}. 2015.
  \doi{10.5281/zenodo.1179130}
  
\item That Syncing Feeling: Networked Strategies for Enabling
    Ensemble Creativity in iPad Musicians (C. Martin, H. Gardner). In
  \emph{Proceedings of CreateWorld}. 2015. \hdl{1885/95216}

\item Preserving Musical Performances on Touch-Screens (C. Martin, H.
  Gardner). In \emph{Proceedings of the CHI 2014 Workshop on Curating
    the Digital: Spaces for Art and Interaction}. 2014.
  \doi{10.5281/zenodo.1175599}

\item Exploring Percussive Gesture on iPads with Ensemble
    Metatone (C Martin, H Gardner, \& B Swift). In
  \emph{Proceedings of the ACM Conference on Human Factors in
    Computing Systems (CHI)}. 2014. \doi{10.1145/2556288.2557226}

\item MetaTravels and MetaLonsdale: iPad Apps for Percussive
    Improvisation (C Martin, H Gardner, \& B Swift). In \emph{CHI
    Extended Abstracts}. 2014. \doi{10.1145/2559206.2574805}

\item Making Improvised Music for iPad and Percussion with
    Ensemble Metatone. In \emph{Proceedings of the Australasian
    Computer Music Conference (ACMC)}. 2014. \hdl{1885/95314}

\item Integrating Mobile Music with Percussion Performance
    Practice (C Martin). In \emph{Proceedings of the International
    Computer Music Conference (ICMC)}. 2013.
  \hdl{2027/spo.bbp2372.2013.073}

\item Performing with a Mobile Computer System for Vibraphone (C
  Martin). In \emph{Proceedings of the International Conference on New
    Interfaces for Musical Expression (NIME)}. 2013.
  \arxiv{2012.00265}

\item Creating Mobile Computer Music for Percussionists: Snow
    Music (C Martin). In \emph{Interactive: Proceedings of the
    Australasian Computer Music Conference (ACMC)}. 2012.
  \doi{10.13140/RG.2.1.5150.5687}

\item Strike on Stage: a Percussion and Media Performance (C
  Martin, \& CH Lai. In \emph{Proceedings of the International
  Conference on New Interfaces for Musical Expression (NIME)}. 2011. \doi{10.5281/zenodo.1178103}

\item Audience Interactive Performance in ``The Last Man to
    Die'' (C Martin, H Cormick, \& B Forster). In
  \emph{Engagement: Proceedings of the Australasian Computer Music
    Conference (ACMC)}. 2010. \hdl{1885/101945}

\item Cross-Artform Performance using Networked Interfaces: Last
    Man to Die's Vital LMTD (C Martin, H Cormick, \& B Forster).
  In \emph{Proceedings of the International Conference on New
  Interfaces for Musical Expression (NIME)}. 2010. \doi{10.5281/zenodo.1177843}

\end{enumerate}

\subsubsection*{Online Press}

\begin{itemize}

\item Shape-shifting robots in the wild: the DyRET robot can rearrange
  its body to walk in new environments (D Howard and C Martin).
  \emph{The Conversation}. 2021.

\end{itemize}

\subsubsection*{Theses}
\begin{itemize}

\item \emph{Apps, Agents, and Improvisation: Ensemble Interaction with
    Touch-Screen Digital Musical Instruments}. PhD Thesis, The
  Australian National University. 2016. \hdl{1885/101786}

\item \emph{Mobile Computer Music for Percussionists}. Master
  Thesis, Lule{\aa} University of Technology. 2012.
  \doi{10.13140/RG.2.1.1688.9681} 

\item \emph{Percussion and Computer in Live Performance}. Master
  Thesis, Australian National University. 2009. \hdl{1885/13813}

\item \emph{Mathematical Foundations of Computer Based Proof
    Verification}. Honours Thesis, Australian National University.
  2008.
\end{itemize}

\pagebreak
\noindent
\includegraphics[width=\textwidth]{header-teaching}


\subsection*{Key Teaching Achievements}
\label{keyteachingachievements}

\begin{itemize}

  \item Led tutor (casual/sessional academic) training at the College of Systems and Society in 2025 and 2026 training $\approx 90$ tutors in 2025 expanding to $\approx 130$ in the first half of 2026.

  \item Initiated Sound and Music Computing course at ANU growing from 7 students in 2019 to 180 in 2025.
  
  \item Co-led design and implementation of Human Centred and Creative Computing major, minor and specialisations at ANU available from 2023, serving as convenor for these sub-plans from 2024.

  \item Led teaching of COMP2300/6300 Computer Organisation and Program Execution in 2019, 2020, and 2022, winning college teaching awards through COVID and significantly redesigning the course in 2022 achieving high SELT results.

\end{itemize}

\subsection*{Teaching Experience}
\label{teachingexperience}

\begin{itemize}

\item \textbf{2025-} Course Convenor for COMP3900/6390 Human-Computer Interaction ($\approx$ 300 students)

\item \textbf{2023-} Course Convenor for COMP4350/8350 Sound and Music Computing. ($\approx$ 180 students)

\item \textbf{2020--2022} Course Convenor for COMP1720/6720, Art and
  Interaction in New Media. ($\approx$ 200 students)
    
\item \textbf{2019--2020, 2022} Course Convenor for COMP2300/6300, Computer Organisation and Program Execution. ($\approx$ 400 students)

\item \textbf{2019--2020, 2022} Course Convenor for COMP2710 Laptop Ensemble
  
\item \textbf{2018--2019} Workshop leader for ``Creative Prediction: with Neural Networks'' at the ALife 2018 and NIME 2019 conferences.

\item \textbf{2018} Co-Lecturer in Advanced Topics in Artificial Intelligence for Intelligent Systems at University of Oslo.

\item \textbf{2016--2019} Masters supervision at the Robotics and Intelligent Systems group, University of Oslo.
  
\item \textbf{2013--2016} Casual Lecturer and Tutor in Percussion and Music Technology at the ANU School of Music.

\item \textbf{2014--2016} Tutor in undergraduate and masters computing courses at the ANU Research School of Computer Science.

\item \textbf{2015} Directed the ``Experimental Music Studio Summer Camp'' and ``DIY Synth Workshop'' at the ANU School of Music

\item \textbf{2014--2015} Coordinated ``Hour of Code'' workshops using the Processing programming environment at the ANU College of Engineering and Computer Science

\item \textbf{2007} Tutor for ``Mathematics and Applications 1 \& 2'', ANU Mathematical Sciences Institute.

\item \textbf{2005--2016} Private Percussion Teacher (including at Canberra Grammar School, Burgmann Anglican School, ANU School of Music)

\end{itemize}

\subsection*{Code Repositories and Released Software}

\begin{itemize}

\item \textbf{Keras Mixture Density Network Layer:} A mixture density
  network (MDN) Layer for Keras using TensorFlow's distributions
  module. \doi{10.5281/zenodo.3526753}

\item \textbf{myo-to-osc:} A pure-python cross-platform solution for
  simply connecting Myo armbands to OSC-connected software. \url{https://github.com/cpmpercussion/myo-to-osc}

\item \textbf{RoboJam:} A Mixture Density RNN and web app for creating
  and responding to musical touchscreen performances.
  \doi{10.5281/zenodo.1064014}

\item \textbf{Gesture-RNN:} Deep recurrent neural network for
  generating ensemble performances of high-level touch gestures.
  (\href{https://github.com/cpmpercussion/gesture-rnn}{github link})

\item \textbf{MicroJam:} Social music-making app for iOS written in
  Swift. (\href{https://github.com/cpmpercussion/microjam}{github
    link})

\item \textbf{PhaseRings:} Musical performance app for iOS written in
  Objective C.
  (\href{https://github.com/cpmpercussion/phaserings}{github link})

\item \textbf{Metatone Classifier:} Ensemble director agent web
  application written in Python with Tornado and Scikit-Learn.
  (\href{https://github.com/cpmpercussion/MetatoneClassifier}{github
    link})

\end{itemize}

\subsection*{Selected Speaking and Media Appearances}

\begin{itemize}
\item \textbf{2021} linux.conf.au (online)
\item \textbf{2020} Data Science Conference Europe (online)
\item \textbf{2019} WIN News \& Nine News, Canberra: Electronic
  Carillon Christmas Installation, Canberra Centre
\item \textbf{2019} Canberra Times, Sherryn Groch: "ANU School of
  Music's new orchestra is far from conventional"
\item \textbf{2018} Tekna Big Data Network, Norway
\item \textbf{2017} Cutting Edge Festival, Oslo Innovation Week
\item \textbf{2017} Technology and Emotions Conference, Oslo
\item \textbf{2016} Music and Video Games Event, Oslo Konserthus
\item \textbf{2014} Percussive Arts Society International Convention,
  Clinic/Performance, Indianapolis, USA
\item \textbf{2014} ANU TV "Colour Music"
\item \textbf{2013} Canberra Times: "It Lives! Theremin returns from dead"
\end{itemize}

\pagebreak
\noindent
\includegraphics[width=\textwidth]{header-cuttingedge}


\subsection*{Grants and Awards}
\label{grantsandawards}

\begin{itemize}

\item \textbf{2022} Australian Government, Defence Science and Technology Group, Artificial Intelligence for Decision Making Initiative (\$30,000)

\item \textbf{2020} Dean's Award for Remote Teaching and Student Experience
  (COMP2300/COMP6300 Computer Organisation and Program Execution).
  
\item \textbf{2019} NVIDIA Corporation Data Science GPU Grant (hardware valued at
  US\$3000).

\item \textbf{2019} Data61 Collaborative Research Project “The Augmented Urban Web”
  (H. Gardner, C. P. Martin, B. Swift, C. Nguyen, M. Adcock, S.
  Anderson, S. Malnar) (\$20,000).

\item \textbf{2017} University of Oslo, Department of Informatics internal
  funding for GPU workstations (\$22,500)
  
\item \textbf{2017} Arts Council Norway (Kulturrådet), project development grant for “Synesthetic” for Marimba, wearable electronics and lighting (C. Hopgood, G. J. Grétarsson, C. P. Martin) (\$11,600)

\item \textbf{2014} ANU Vice Chancellor's HDR Travel Grant to present at Percussive Arts Society International Convention (\$1,000)

\item \textbf{2013} Australian Postgraduate Award (APA) and ANU Supplementary Scholarship 

\item \textbf{2011} National Science Week Seed Grant to produce “The Last Man to Die” performance in Canberra. (H. Cormick, B. Forster, C. P. Martin) (\$1482)

\item \textbf{2011} ArtsACT Out of Round Funding to attend University of Delaware Jazz Vibraphone Workshop (\$1,500)

\item \textbf{2010} Australia Council for the Arts, Theatre Section, New Work - Young Artists Initiative to produce “The Last Man to Die” at the Blue Room Theatre, Perth (H. Cormick, B. Forster, C. P. Martin) (\$14,400)

\item \textbf{2010} Department of Innovation, Industry, Science and Research, National Enabling Technologies Strategy Funding to develop “The Last Man to Die” production (H. Cormick, B. Forster, C. P. Martin) (\$9,800)

\item \textbf{2010} ArtsACT Project Funding to develop “Strike on Stage” performance for percussion and interactive media (C.-H. Lai, C. P. Martin) (\$10,840)

\item \textbf{2010} ArtsACT Project Funding to develop “The Last Man to Die” (C. P. Martin, H. Cormick, B. Forster) (\$9,000)

%\item \textbf{2009} Space Program Residency, PACT Centre for Emerging Artists, Sydney (C. P. Martin, H. Cormick, B. Forster) (in kind support for theatre hire and artistic development)
%\item \textbf{2009} Emerging Artists Residency, Canberra Youth Theatre (H. Cormick, B. Forster, C. P. Martin) (in kind support for theatre hire and artistic development)
%\item \textbf{2006} ANU Research School of Information Sciences and Engineering Summer Scholarship

\end{itemize}

\pagebreak
\noindent
\includegraphics[width=\textwidth]{header-albums}

\subsection*{Selected Audio Releases}
\label{selectedrecordings}

\begin{itemize}

\item \textbf{2026} \emph{Hyphae - Charles Martin and Alexander Hunter}. An album of improvised music for electronic duo with AI agents blending synthesised and sampled sounds in an unusual and unexpected dialogue with machine learning algorithms (in production)

\item \textbf{2023} \emph{Orbits - Charles Martin} Performed live on iPad and synthesisers, 2023-08-03 at the Powerhouse Museum, Ultimo, Sydney. This album was a live soundtrack composed for a magic lantern show by Martyn Jolly with astronomical and terrestrial slides from the 1880s-1900s and contemporary slides by Anna Madeleine Raupach.
(\href{https://charlesmartin.bandcamp.com/album/orbits}{Bandcamp link})

\item \textbf{2017} \emph{Andromeda is Coming.} Andromeda is Coming
  (Charles Martin and Alexander Hunter) are a Canberra-based duo
  utilising viola da gamba, percussion and live electronics. This
  debut album collects a series of free-improvised sections across two
  field recordings, each capturing Canberra's intersection of urban
  and bush environments. (\href{http://collectedresonances.bandcamp.com/album/andromeda-is-coming}{Bandcamp link})

\item \textbf{2014} \emph{Ensemble Metatone.} Improvised works from
  Ensemble Metatone, a group dedicated to creating experimental works
  with new iPad instruments, these works were performed with a new
  vocabulary of touch-gestures inspired by their percussive
  backgrounds.
  (\href{http://charlesmartin.bandcamp.com/album/ensemble-metatone}{Bandcamp
    link})

\item \textbf{2013} \emph{Nordlig Vinter - Charles Martin.}
  Compositions for percussion and iOS devices created by Charles
  Martin while living in Pite{\aa}, near the Arctic circle in northern
  Sweden. Harshly cut field recordings reflect the dry, icy
  environment while minimal jazz inspired melodies on vibraphone and
  marimba evoke the quiet beauty of such a strange place.
  (\href{http://charlesmartin.bandcamp.com/album/nordlig-vinter}{Bandcamp
    link})

\item \textbf{2012} \emph{Sounds from the Treetops - Ensemble
    Evolution.} Six works for percussion trio, voice, and clarinet,
  composed and performed by Ensemble Evolution in collaboration with
  the iconic Treehotel in Northern Sweden.
  (\href{http://ensembleevolution.bandcamp.com/album/sounds-from-the-treetops}{Bandcamp
    link})

\item \textbf{2012} \emph{Solstormen - Ensemble Evolution.} A suite of
  mallet trios by the members of Ensemble Evolution drawing elements
  from jazz, minimalism and contemporary classical music. This music
  recalls the aurora borealis and the eerie landscape of Northern
  Sweden.
  (\href{http://ensembleevolution.bandcamp.com/album/solstormen-live}{Bandcamp
    link})
\end{itemize}

\pagebreak
\noindent
\includegraphics[width=\textwidth]{header-metatone}


\subsection*{Selected performances and productions}
\label{selectedperformancesandproductions}

\begin{itemize}

\item \textbf{2026} \emph{SoundOut Fesitval} performances at Drill Hall Gallery, February 2026. \url{https://charlesmartin.au/blog/2026/02/01/soundout-2026}

\item \textbf{2025} \emph{SMCClablive \#1} performance at ANU School of Music, November 2025. \url{https://smcclab.au/2025/11/24/smcclablive-wrapup.html}

\item \textbf{2024} \emph{AI Musical Instrument Free Improvised Music Practice} performances at Drill Hall Gallery, Canberra Contemporary Art Space Manuka, Share house garage in O'Connor.

\item \textbf{2024} \emph{Andromeda is Coming and Artificially Intelligent Friends} performance at ANU School of Music, May 2024. \url{https://smcclab.github.io/2024/05/16/andromeda-concert.html} Video: \url{https://youtu.be/vH4qWDztFNk}

\item \textbf{2022-23} \emph{Cubing Sound}, with Yichen Wang (ANU School of Computing). AR and touchscreen performances at OzCHI2022 and NIME2023 (e.g., \url{https://youtu.be/uR7LbpRMMVM}).

\item \textbf{2020-21} \emph{Heritage in the Limelight: The Magic
    Lantern in Australia and the World}, musician for Magic Lantern
  performances led by Dr Elisa deCourcy and Prof Martyn Jolley (ANU
  School of Art and Design) (e.g., \url{https://vimeo.com/405336569})

% \item \textbf{2020} \emph{Sci-Fi Theremin Jukebox Performance}, with
%   Kit Devine (ANU School of Art and Design), Alec Hunter (ANU School
%   of Music), for the ANU Centre for Public Awareness of
%   Science ZOOMPOSIUM: Performance, Science \& Technology.

\item \textbf{2020} \emph{Neurofeedback - solo performance for
    computer and artificial neural network} performed at \emph{Smoke}
  (ANU School of Music, Jan 2020) (\url{https://youtu.be/WumtMGHAuV8})

\item \textbf{2019-2020} \emph{ANU Laptop Ensemble}, co-director of
  ANU LENS since 2019, involved as performer and mentor for student
  performances in various venues and online (e.g.,
  \url{https://youtu.be/24YMz90b6bo})

\item \textbf{2018 jun} \emph{Stillness Under Tension}, ensemble performance
  for Myo muscle sensors and single board computers at NIME 2018,
  Virginia Tech, USA.

\item \textbf{2017 aug} \emph{Metatone iPad Ensemble}, 3-hour
  improvised touchscreen performance at ``Elvelangs i Fakkellys'',
  Oslo. \url{https://youtu.be/NepiJe-TB_Q}

\item \textbf{2017 jul} \emph{Andomeda is Coming}, album launch at
  Gorman House Arts Centre, Canberra.
  \url{https://youtu.be/QfLQb5dPWfY}

\item \textbf{2017 jun} \emph{Synesthetic}, prototype performance of
  works for marimba and computer controlled lighting at Nydans, Oslo.
  \url{https://youtu.be/F9IcC7vv4QQ}

% \item \textbf{2017 may} \emph{Music Tech p{\aa} IFI}, showcase concert
%  at UiO Department of Informatics. 

% \item \textbf{2016 jul} \emph{PhaseRings for iPad Ensemble and
%     Ensemble Director Agent}, invited performance at the International
%   Conference on Auditory Display, Llewellyn Hall, Canberra.
%   \url{https://youtu.be/aDEQMLwd8ok}

% \item \textbf{2015 nov} \emph{Andromeda is Coming}, Charles Martin and Alexander Hunter, experimental music performance at ElectroFringe 2015, 107 Projects, Sydney.

%\item \textbf{2014 nov} \emph{iPads in Percussion Ensemble}, seminar and performance at the 2014 Percussive Arts Society International Convention, Indianapolis, USA.

% \item \textbf{2014 aug} \emph{Ensemble Metatone at Colour Music}, iPad/Percussion ensemble performance as part of the ANU Drill Hall Gallery's \emph{Colour Music} Exhibition. \url{https://youtu.be/70xlD1zrfsk}

% \item \textbf{2014 may} \emph{Frost Beat}, Frost Ice Bar, Boston, USA. Improvised performance series with percussion and ice instruments. \url{https://youtu.be/Sis3XmSzafs}

% \item \textbf{2014 mar} \emph{Loner Series: Ensemble Metatone}, You
% Are Here Festival, Canberra, Performance of new works for iPad and %
% percussion ensemble. \url{https://youtu.be/F9IcC7vv4QQ}

% \item \textbf{2013 nov} \emph{Ensemble Evolution: Sounds from the Treetops}, ensemble showcase concert at the 2013 Percussive Arts Society International Convention, Indianapolis, USA.

%\item \textbf{2012 jun} \emph{Sounds from the Treetops}, Commission with \emph{\href{http://www.ensemble-evolution.com}{Ensemble Evolution}} for \href{http://www.treehotel.se}{Treehotel}, an iconic boutique hotel in Northern Sweden. (\href{http://charlesmartin.com.au/blog/2012/10/8/sounds-from-the-treetops-ensemble-evolution}{blog and video})

%\item \textbf{2012 jun} \emph{Ensemble Evolution - NattMusik Tour},
% performing new music for percussion throughout Northern Sweden in
% celebration of the midnight sun.
% (\href{http://charlesmartin.com.au/2012/10/8/nattmusik-with-ensemble-evolution}{blog
% and videos})

% \item \textbf{2012 apr} \emph{Sikt} w\slash Malin H\"{u}lphers, Anders
% {\AA}strand, Sanna Hodell, Fredrik Strid. Stockholm, Sweden.
% Performing live electronic sound in Stockholm.
% (\href{http://charlesmartin.com.au/blog/2012/4/19/sikt-in-stockholm.html}{blog
% article})

%\item \textbf{2012 mar} \emph{Ensemble Evolution - Solstormen} tour, Australia. Featuring performances and workshops using \emph{Snow Music}, a percussion work and computer instrument for iOS in Melbourne, Sydney and Canberra. \url{https://youtu.be/Ud3Mq_LLMGI}

% \item \textbf{2011 nov} Workshop and duo concert, \emph{drums +
% gadgets}, with Noah Demland in Columbus Ohio. Featuring new mobile
% computer instruments for iOS devices.
% (\href{http://charlesmartin.com.au/blog/2012/1/17/columbus-collaboration.html}{blog
% article})

% \item \textbf{2011 sep} \emph{Norra Vinter} at Electrofringe, duo
% performance with Christina Hopgood of new works for percussion and iOS
% devices.
% (\href{http://charlesmartin.com.au/blog/2011/9/29/nordlig-vinter-at-electrofringe.html}{blog
% article})

% \item \textbf{2011 mar}
% \emph{\href{http://www.ensemble-evolution.com}{Ensemble Evolution}} at
% Pite{\aa} Percussion Repertoire Festival, Studio Acusticum, Pite{\aa}
% and Kroumata's Capitol Theatre, Stockholm. Including ensemble and solo
% works by \emph{Ensemble Evolution} and other emerging composers.
% (\href{http://www.youtube.com/watch?v=ltmNTijFulo}{video})

% \item \textbf{2010 nov} \emph{Swexican Trio} w\slash Anders
% {\AA}strand, (performing live electronic sound + video), Percussive
% Arts Society International Convention, Indianapolis, USA.

%\item \textbf{2010 oct} \emph{Anders {\AA}strand, Dave Samuels,
%Evaristo Aguilar w\slash Ensemble Evolution} - Studio Acusticum,
%Pite{\aa}, Kroumata's Capitol Theatre, Stockholm.

%\item \textbf{2010 feb-oct} \emph{``THE LAST MAN TO DIE''} - Blue Room Theatre, Perth; Brisbane Festival, Under the Radar programme, Metro Arts, Brisbane; Street 2, The Street Theatre, Canberra; PACT Centre for Emerging Artists, Sydney; Belconnen Theatre, Canberra. (\href{http://www.lastmantodie.net/p/show.html}{website})

% %\item \textbf{2010 jul} \emph{Strike on Stage Micro} and
% \emph{Improvisation for Touch Table in Supercollider} - CCRMA
% Listening Room, Stanford University, USA.
% (\href{http://charlesmartin.com.au/2010/8/15/stanford-ccrma-tabletop-interface-workshop.html}{blog
% article and video})

% \item \textbf{2010 jun} \emph{Strike on Stage: Percussion and Media Performance}, Belconnen Arts Centre, Canberra, New Interfaces for Musical Expression Conference, ABC Eugene Goosens Hall, Sydney. (\href{http://charlesmartin.com.au/strike-on-stage/}{blog and videos})

%\item \textbf{2009 sep-dec} \emph{Vital LMTD} with \emph{Last Man to
% Die} cross art-form group - This Not Art festival, Newcastle, Street
% 2, The Street Theatre, Canberra, Canberra Youth Theatre.
% (\href{http://charlesmartin.com.au/blog/2009/11/3/vital-lmtd-at-the-street-theatre.html}{blog
% article})

% \item \textbf{2008 sep-nov} \emph{Cognition Series}, with \emph{The
% Hunting Season} cross-artform group - Belconnen Theatre, Street 2, The
% Street Theatre, Canberra and La Mama Theatre, Melbourne.
% (\href{http://charlesmartin.com.au/blog/2008/9/21/cognition.html}{blog
% article})

\end{itemize}

% \pagebreak

% \subsection*{References}

% \subsubsection*{Research References}
% \begin{itemize}

% \item \textbf{Prof Roger Dannenberg, Professor of Computer Science,
%     Art \& Music, Carnegie Mellon University} \\
%   T: 412-268-3827 \\
%   E: rbd@cs.cmu.edu

% \item \textbf{Dr Rebecca Fiebrink, Reader, Creative Computing
%     Institute, University of the Arts, London} \\
%   E: r.fiebrink@arts.ac.uk

% \item \textbf{Prof Roger Dean, Research Professor of Sonic
%     Communication, The MARCS Institute for Brain, Behaviour and
%     Development, Western
%     Sydney University.}\\
%   E: roger.dean@westernsydney.edu.au \\
%   T: +61 2 9772 6902

% \item \textbf{Ge Wang, Associate Professor of Music and Computer Science, Centre for Computer
%     Research in Music and Acoustics (CCRMA), Stanford University} \\
%   E: ge@ccrma.stanford.edu \\
%   (N.B.: may be unavailable over USA summer until September 2021)
% \end{itemize}

% \subsubsection*{Supervisors and Colleagues}
% \begin{itemize}
% \item \textbf{Prof Jim T{\o}rresen, Professor, Department of
%     Informatics, University of Oslo}\\
%   T: (+47) 2285 2454 \\
%   E: jimtoer@ifi.uio.no \\
%   A: Ole Johan Dahls hus, Gaustadal{\'e}en 23 B, N-0373 OSLO, Norge

% \item \textbf{Dr Kyrre Glette, Associate Professor, Department of
%     Informatics, University of Oslo} \\
%   T: (+47) 2284 1695 \\
%   E: kyrrehg@ifi.uio.no \\
%   A: Ole Johan Dahls hus, Gaustadal{\'e}en 23 B, N-0373 OSLO, Norge

% \item \textbf{Prof Tony Hosking, Professor, ANU School of
%     Computing}  \\
%   T: (+612) 6125 9358 \\
%   E: Tony.Hosking@anu.edu.au \\
%   A: CSIT Building 108, The Australian National University, Canberra,
%   ACT, 2601, Australia

% \item \textbf{Dr Henry Gardner, Honorary Associate Professor, ANU School of Computing} \\
%   T: (+612) 6125 8181 \\
%   E: Henry.Gardner@anu.edu.au \\
%   A: Hanna Neumann Building 145, The Australian National University,
%   Canberra, ACT, 2601, Australia

% \end{itemize}

\end{document}  
